% Kesimpulan
\section{Kesimpulan}

Paper ini mengestimasi struktur korelasi kedelapan subtes UTBK-SNBT 2026 dengan memanfaatkan data \textit{crowdsourced} dari platform X ($N=1.886$), yang menghasilkan rerata korelasi \textit{off-diagonal} pada \textit{raw sample} sebesar 0,300. Tiga indikator utama mendukung kesimpulan bahwa nilai 0,300 berfungsi sebagai batas bawah dari korelasi populasi riil akibat \textit{restriction of range}: (1) deviasi parameter sampel terhadap parameter nasional tidak resmi, (2) dinamika fluktuasi dan pelandaian korelasi sepanjang tahap penjaringan data, serta (3) peningkatan rerata korelasi dari 0,300 menjadi 0,425 setelah koreksi Thorndike Case III.

Koreksi Thorndike Case III berbasis iterasi titik tetap matriks kovarians penuh pada \textit{pipeline} komputasi terbuka R menghasilkan estimasi korelasi populasi terkoreksi sebesar $\bar{R} \approx 0,425$ ($S_Z \approx 67,72$). Interval \textit{bootstrap} 95\% untuk $\bar{R}$ ($[0,407; 0,442]$) hanya menangkap ketidakpastian pengambilan sampel. Ketergantungan pada simpangan baku ($SD$) nasional tidak resmi jauh lebih besar, sebab perubahan seluruh $SD$ sebesar $\pm 10\%$ menggeser $\bar{R}$ menjadi $0,300$--$0,519$. Dengan demikian, estimasi ini bersifat kondisional terhadap parameter nasional yang belum terverifikasi. Nilai itu menunjukkan bahwa korelasi populasi diperkirakan lebih tinggi daripada korelasi \textit{raw sample}. Matriks translasi persentil standar yang dihasilkannya (lihat Hasil) lebih tepat dipakai sebagai perkiraan kasar posisi persentil nasional dari skor komposit tertentu.

Analisis lanjutan melalui CFA menunjukkan bahwa struktur laten kedelapan subtes lebih baik direpresentasikan oleh model dua faktor (Literasi dan Numerasi, $\phi = 0,813$) daripada model dengan \textit{g-factor} tunggal. Analisis itu juga mengidentifikasi variasi homogenitas korelasi antarsubtes yang membatasi presisi koreksi Case III pada tingkat pasangan subtes individual. Perlakuan terhadap dua skor LBI, yang berasal dari satu tes, tidak mengubah kesimpulan utama, sebab $\bar{R}$ berkisar $0,414$--$0,425$ ketika salah satu atau kedua skor LBI dikeluarkan.

Statistik nasional tidak resmi menunjukkan distribusi skor yang miring ke kanan dengan ekor kiri terpotong. Simulasi dengan marginal Johnson SB menunjukkan bahwa asumsi normal berdampak kecil di sekitar $\pm 1\,SD$ tetapi besar di ekor, misalnya skor komposit pada persentil normal ke-99,9 sekitar $730$ pada simulasi dibandingkan $703$ pada normal. Bila korelasi antarsubtes menurun pada skor tinggi (SLODR), koreksi Case III menaksir rendah korelasi populasi sehingga $\bar{R} \approx 0,425$ lebih tepat dibaca sebagai batas bawah, dan efek itu sebagian meniadakan efek marginal miring di ekor atas. Data paper ini tidak dapat membedakan SLODR dari \textit{restriction of range}.

Paper ini memberikan kontribusi rangkap tiga: \textit{kontribusi substantif}, berupa estimasi empiris awal dan matriks terkoreksi antarsubtes UTBK-SNBT 2026 yang hingga kini belum tersedia dalam dokumen publik, lengkap dengan analisis ketidakpastian dan sensitivitasnya; \textit{kontribusi struktural}, berupa konfirmasi CFA atas struktur dua faktor Literasi dan Numerasi yang mendasari kedelapan subtes, sekaligus diagnostik homogenitas korelasi yang menjadi dasar kehati-hatian dalam menafsirkan koreksi Case III; serta \textit{kontribusi metodologis}, berupa temuan bahwa bias seleksi pada data \textit{crowdsourced} media sosial tampaknya bersumber dari saringan sosiologis audiens platform yang tidak dapat dipulihkan semata-mata oleh formulasi \textit{framing} pengumpulan data. Estimasi korelasi dan persentil yang definitif pada akhirnya menuntut keterbukaan akses terhadap data mikro tingkat populasi dari Panitia SNPMB.